\documentclass[11pt,a4paper]{article}
\usepackage{amsmath,amssymb}
\usepackage{listings}
\usepackage{hyperref}
\usepackage[margin=1in]{geometry}
\usepackage{fancyhdr}

\pagestyle{fancy}
\fancyhf{}
\fancyhead[L]{Module 08: Structures and Unions}
\fancyhead[R]{C Programming Course}
\fancyfoot[C]{\thepage}

\lstset{
  language=C,
  basicstyle=\ttfamily\small,
  keywordstyle=\bfseries,
  commentstyle=\itshape,
  numbers=left,
  numberstyle=\tiny,
  frame=single,
  breaklines=true,
  tabsize=4
}

\title{Module 08: Structures and Unions}
\author{C Programming Course}
\date{}

\begin{document}
\maketitle
\thispagestyle{fancy}

\section{Defining Structures}
\begin{lstlisting}
#include <stdio.h>

struct Point {
    double x;
    double y;
};

int main(void) {
    struct Point p = {3.0, 4.0};
    printf("(%.1f, %.1f)\n", p.x, p.y);
    return 0;
}
\end{lstlisting}

\section{Nested Structures}
\begin{lstlisting}
struct Address {
    char city[50];
    int zip;
};

struct Person {
    char name[50];
    struct Address addr;
};
\end{lstlisting}

\section{Arrays of Structures}
\begin{lstlisting}
struct Student {
    char name[30];
    int grade;
};

struct Student class[3] = {
    {"Alice", 90}, {"Bob", 85}, {"Carol", 92}
};
\end{lstlisting}

\section{Pointers to Structures}
Use the arrow operator (\texttt{->}) to access members through a pointer.

\begin{lstlisting}
#include <stdio.h>

struct Point {
    double x;
    double y;
};

void print_point(const struct Point *p) {
    printf("(%.1f, %.1f)\n", p->x, p->y);
}

int main(void) {
    struct Point pt = {1.0, 2.0};
    print_point(&pt);
    return 0;
}
\end{lstlisting}

\section{Unions and Bit Fields}
A union shares memory among its members; only one member is valid at a time.

\begin{lstlisting}
union Data {
    int i;
    float f;
    char c;
};
\end{lstlisting}

Bit fields allow packing flags into minimal space:
\begin{lstlisting}
struct Flags {
    unsigned int active : 1;
    unsigned int ready  : 1;
    unsigned int mode   : 3;
};
\end{lstlisting}

\section{typedef and enum}
\begin{lstlisting}
typedef struct {
    double x, y;
} Vec2;

typedef enum { RED, GREEN, BLUE } Color;
\end{lstlisting}

\texttt{typedef} creates an alias for a type; \texttt{enum} defines named integer constants.

\end{document}
